% Shared pieces for the lr0_mosfet cross-section cartoons.
%
% Two families use this file:
%  - the field-effect intro (mosfet_off/subthreshold/strong_inversion...):
%    white background, glowing blue n+ blocks, gate box, S/G/D terminals
%  - the carrier cartoons (accumulated/depleted/weakinv/vds_*): gray
%    substrate, n+ boxes with blue minus signs, red plus signs in the
%    p- bulk, teal oxide strip under the gate
%
% The colour language follows the original hand drawings: free electrons
% are blue, holes are red, the lecture prose ("imagine the free electrons
% as a blue fluorescent color") depends on it.

\definecolor{echarge}{RGB}{30,60,220}   % electron blue
\definecolor{hcharge}{RGB}{220,30,30}   % hole red
\definecolor{oxteal}{RGB}{60,140,160}   % oxide/channel strip

% ------------------------------------------------------------------
% Family 1: field-effect intro frame.
% Geometry: n+ blocks [0,2.9] and [5.9,8.8] x [-1.5,0], gate [2.6,6.2].
% ------------------------------------------------------------------
\newcommand{\fetblock}[2]{% x0 x1: glowing blue block below the surface
  \fill[echarge!30, rounded corners=8pt] (#1-0.12,-1.62) rectangle (#2+0.12,0.10);
  \fill[echarge, rounded corners=6pt] (#1,-1.5) rectangle (#2,0);
}

\newcommand{\fetterminal}[3]{% x y label: stem with a small open circle
  \draw[thick] (#1,#2) -- (#1,#2+0.35);
  \draw[thick] (#1,#2+0.44) circle (0.09);
  \node[anchor=south] at (#1,#2+0.62) {#3};
}

\newcommand{\fetframe}{%
  \fetblock{0}{2.9}
  \fetblock{5.9}{8.8}
  % gate: rounded box sitting just above the surface
  \draw[line width=1.6pt, rounded corners=10pt] (2.6,0.12) rectangle (6.2,1.05);
  \node[anchor=center] at (4.4,0.58) {GATE};
  \fetterminal{1.3}{0}{S}
  \fetterminal{4.4}{1.05}{G}
  \fetterminal{7.4}{0}{D}
}

% ------------------------------------------------------------------
% Family 2: carrier cartoon frame.
% Surface at y=0. Substrate down to y=-3.6, x in [-0.4,10.4].
% n+ boxes [0.3,2.9] and [7.1,9.7] x [-1.15,0]. Gate [3.0,7.0].
% ------------------------------------------------------------------
\newcommand{\eminus}[2]{\node[echarge, anchor=center, scale=0.7] at (#1,#2) {$\boldsymbol{-}$};}
\newcommand{\hplus}[2]{\node[hcharge, anchor=center, scale=0.8] at (#1,#2) {$\boldsymbol{+}$};}

% rows of minus signs filling an n+ box
\newcommand{\efill}[2]{% x0 x1
  \foreach \y in {-0.25,-0.55,-0.85}{
    \foreach \x in {0.15,0.55,...,2.4}{
      \eminus{#1+\x}{\y}
    }
  }
}

\newcommand{\bulkframe}{%
  % substrate outline
  \draw[line width=2.2pt, gray!60, rounded corners=4pt]
    (-0.4,0) -- (10.4,0) -- (10.4,-3.6) -- (-0.4,-3.6) -- cycle;
  % n+ wells
  \draw[line width=1.4pt, rounded corners=3pt] (0.3,0) -- (0.3,-1.15) -- (2.9,-1.15) -- (2.9,0);
  \draw[line width=1.4pt, rounded corners=3pt] (7.1,0) -- (7.1,-1.15) -- (9.7,-1.15) -- (9.7,0);
  \efill{0.3}{2.9}
  \efill{7.1}{9.7}
  \node[anchor=south west, scale=0.8] at (0.32,-1.12) {$n+$};
  \node[anchor=south east, scale=0.8] at (9.68,-1.12) {$n+$};
  % deep bulk holes
  \foreach \y/\xoff in {-3.3/0.0, -2.85/0.25}{
    \foreach \x in {0.1,0.7,...,10.2}{
      \hplus{\x+\xoff}{\y}
    }
  }
  \node[hcharge!70!black, anchor=center] at (5.0,-3.95) {$p-$};
  % oxide strip and gate box
  \fill[oxteal!80] (3.0,0.02) rectangle (7.0,0.24);
  \draw[line width=1.6pt, rounded corners=6pt] (2.75,0.26) rectangle (7.25,1.15);
}

% one row of holes between the wells
\newcommand{\neckrow}[2]{% y xoff
  \foreach \x in {3.2,3.8,...,6.9}{
    \hplus{\x+#2}{#1}
  }
}

% holes to the sides underneath the wells
\newcommand{\sideholes}{%
  \foreach \y/\xoff in {-1.6/0.1, -2.15/0.35}{
    \foreach \x in {-0.1,0.5,...,2.6}{ \hplus{\x+\xoff}{\y} }
    \foreach \x in {7.3,7.9,...,10.1}{ \hplus{\x+\xoff}{\y} }
  }
}

% terminals with ground symbols, and the p- label
\newcommand{\bulkterminalgnd}[3]{% x label side(-1 left / 1 right)
  \draw[thick] (#1,0) -- (#1,0.5) -- ({#1+#3*0.7},0.5);
  \draw[thick] ({#1+#3*0.7},0.35) -- ({#1+#3*0.7},0.65);
  \draw[thick] ({#1+#3*0.85},0.42) -- ({#1+#3*0.85},0.58);
  \node[echarge, anchor=south] at (#1,0.65) {#2};
}
